\def\changesize#1{\fontsize{#1}{10.5pt}\selectfont}
\begin{table}[h]
\centering
\setlength{\tabcolsep}{4pt}
\small
\begin{tabular}{cccc}
\toprule
\multicolumn{2}{c}{\textbf{Data Split}} & \multicolumn{2}{c}{\textbf{Average Performance (\%)}} \\
\cmidrule(lr){1-2} \cmidrule(lr){3-4}
\textbf{\sft{}} & \textbf{\opsd{}} & \textbf{\sft{} Model} & \textbf{\ours{} Model} \\
\midrule
6K & 9K & 57.6 & 60.3 \\
9K &  6K& 57.8 & 59.1 \\
7.5K & 7.5K & 57.8 & 59.8 \\
\bottomrule
\end{tabular}
\caption{\textbf{Ablation of data split between \sft{} and \opsd{} phases.} On \qwen{}, we vary how the training data is allocated across the two phases while keeping the total data budget fixed. Allocating more data to \sft{} slightly improves \sft{} model performance, but does not lead to the best final \ours{} accuracy. The best overall result is achieved by assigning more data to \opsd{}, suggesting that once \sft{} has sufficiently elicited self-revision capability, additional data is more effectively used in \opsd{} phase to distill and refine this behavior into stronger generation.}
\label{tab:data-split-ablation}
\end{table}